% =====================================================================
%  1bat-ud1-1.tex  —  HORA 1: Diagnosi inicial i benvinguda a la matèria
% =====================================================================
\horatitol{Sessió 1 — Diagnosi inicial i benvinguda a la matèria}%
{Per què estudiarem matemàtiques aquest curs, i una diagnosi breu i sense nota per saber d'on partim.}

\subsection{Idea clau}

Benvingudes i benvinguts a \textbf{Matemàtiques Aplicades I}. Aquest curs no farem
matemàtiques «abstractes»: partirem sempre de situacions \textbf{reals} de l'àmbit
social i comunitari —ajuts, llindars de renda, edats límit, pressupostos públics— i
hi posarem les eines matemàtiques que ens calguin per entendre-les i decidir-hi.

\begin{idea}
\textbf{Com treballarem}
\begin{itemize}[leftmargin=1.4em, itemsep=2pt]
  \item \textbf{Aules pensants}: molts dies treballareu en grup, drets, a les
        pissarres verticals, abans que us doni cap fórmula. Primer pensar, després
        formalitzar.
  \item \textbf{Contextos reals}: cada bloc arrenca d'un cas autèntic (per exemple, el
        \emph{bo social elèctric}, que veurem aquesta setmana).
  \item \textbf{Pensament crític amb dades}: no només calcular, sinó preguntar-nos si
        un resultat \emph{té sentit}.
\end{itemize}
La diagnosi d'avui \textbf{no té nota}: serveix perquè jo sàpiga d'on partim i us
acompanyi millor.
\end{idea}

\subsection{Exemples}

\begin{exemple}[d'una frase quotidiana a un símbol]
Moltes condicions reals s'amaguen darrere d'una frase. Fixem-nos en com es tradueixen:
\begin{itemize}[leftmargin=1.4em, itemsep=2pt]
  \item «La renda ha de ser \textbf{com a màxim} de $\num{12000}\eur$» $\rightarrow$
        $\text{renda}\le \num{12000}$.
  \item «Cal tenir \textbf{com a mínim} $16$ anys» $\rightarrow$ $\text{edat}\ge 16$.
  \item «El preu \textbf{no pot superar} els $50\eur$» $\rightarrow$ $\text{preu}\le 50$.
\end{itemize}
El símbol $\le$ es llegeix «menor o igual» i el $\ge$, «major o igual». Aquesta
traducció de paraules a símbols és el cor de la unitat.
\end{exemple}

\begin{exemple}[ordenar nombres a la recta]
Per comparar nombres, els situem sobre una \textbf{recta numèrica}: com més a la
dreta, més grans. Ordenem $-3$, \ $1,5$, \ $-\tfrac{1}{2}$, \ $2$:
\[
-3 \;<\; -\tfrac{1}{2} \;<\; 1,5 \;<\; 2.
\]
Compte amb els negatius: $-3$ és \emph{més petit} que $-\tfrac{1}{2}$, encara que «el
$3$» sigui més gran que «el mig».
\end{exemple}

\subsection{Exercicis resolts}

\begin{exresolt}[traduir frases a símbols]
Tradueix cada frase a una desigualtat amb el símbol correcte ($<$, $>$, $\le$, $\ge$):
\begin{enumerate}[label=\alph*), leftmargin=1.6em]
  \item «Per inscriure's cal haver fet més de $18$ anys.»
  \item «L'aforament és com a màxim de $200$ persones.»
\end{enumerate}
\solucio
Identifiquem si l'extrem hi entra o no:
\begin{enumerate}[label=\alph*), leftmargin=1.6em]
  \item «\textbf{més de} $18$» exclou el $18$: $\text{edat}>18$.
  \item «\textbf{com a màxim} $200$» inclou el $200$: $\text{aforament}\le 200$.
\end{enumerate}
\resposta{a) $\text{edat}>18$; \quad b) $\text{aforament}\le 200$.}
\end{exresolt}

\begin{exresolt}[repàs de càlcul bàsic]
Calcula sense calculadora: (a) $\tfrac{1}{2}+\tfrac{1}{3}$; \ (b) el $15\%$ de $80$;
\ (c) resol $2x+1=9$.
\solucio
\textbf{(a)} Reduïm a comú denominador ($6$): $\tfrac{3}{6}+\tfrac{2}{6}=\tfrac{5}{6}$.

\textbf{(b)} El $15\%$ de $80$ és $0,15\cdot 80=12$.

\textbf{(c)} $2x+1=9 \Rightarrow 2x=8 \Rightarrow x=4$.
\resposta{(a) $\tfrac{5}{6}$; \quad (b) $12$; \quad (c) $x=4$.}
\end{exresolt}

\subsection{Exercicis per fer (diagnosi)}

\noindent\textit{Resol-los sols i en silenci. No tenen nota: són per veure què
controlem ja i què haurem de repassar.}

\begin{exercicis}
\begin{exlist}
  \item Calcula: (a) $-7+12$; \quad (b) $-3\cdot(-4)$; \quad (c) $\tfrac{3}{4}-\tfrac{1}{2}$;
        \quad (d) $2,5+1,75$.
  \item Calcula el $20\%$ de $150$.
  \item Resol l'equació $3x-5=10$.
  \item Ordena de menor a major: $\;-2,\ \tfrac{1}{2},\ -\tfrac{5}{2},\ \sqrt{9},\ 0$.
  \item Tria el símbol correcte ($<$, $>$, $\le$ o $\ge$) per a cada frase:
        \begin{enumerate}[label=\alph*), leftmargin=1.6em, topsep=2pt]
          \item «Cal tenir com a mínim $14$ anys.» $\rightarrow$ $\text{edat}\;\rule{1cm}{0.4pt}\;14$.
          \item «La despesa no pot passar de $300\eur$.» $\rightarrow$ $\text{despesa}\;\rule{1cm}{0.4pt}\;300$.
        \end{enumerate}
  \item Calcula la potència $2^5$ i escriu quants zeros té el nombre $10^4$.
  \item \textbf{(Reflexió)} De tota aquesta llista, què et sembla \emph{més difícil}?
        Escriu-ho en una frase. (No hi ha resposta correcta: és per ajudar-me a
        acompanyar-te.)
\end{exlist}
\end{exercicis}

% ---------------------------------------------------------------------
\fullsolucions{Sessió 1}
\begin{solucions}
\begin{sollist}
  \item (a) $5$; \quad (b) $12$; \quad (c) $\tfrac{1}{4}$; \quad (d) $4,25$.
  \item El $20\%$ de $150$ és $0,20\cdot 150=30$.
  \item $3x-5=10 \Rightarrow 3x=15 \Rightarrow x=5$.
  \item $-\tfrac{5}{2}\,(=-2,5) \;<\; -2 \;<\; 0 \;<\; \tfrac{1}{2} \;<\; \sqrt{9}\,(=3)$.
  \item (a) $\text{edat}\ge 14$; \quad (b) $\text{despesa}\le 300$.
  \item $2^5=32$; \quad $10^4=\num{10000}$ té $4$ zeros.
  \item Resposta oberta.
\end{sollist}
\end{solucions}
